\chapter{The Observation-Action Safety Gap and Claim Boundary}
\label{ch:oasg-claim-boundary}
\label{ch:safety-illusion}

The observation-action safety gap (OASG) is the central organizing object of the
monograph.  It names the event in which an action is admissible on the observed state
while unsafe on the true state.  In practical terms, this happens whenever sensing,
transport, synchronization, or estimation degrades more severely than the downstream
controller assumes.  The action remains \emph{observationally legitimate} but becomes
\emph{physically illegitimate}.

\section{A structural hazard rather than a domain bug}
ORIUS treats this as a universal hazard because the structure is independent of plant
details.  Battery dispatch exposes it through hidden state-of-charge violations.  Vehicle
control exposes it through stale headway and closing-speed estimates.  Industrial systems
expose it through delayed thermal or pressure readings.  Healthcare exposes it through
stale alarms and delayed physiological updates.  Navigation and aerospace expose it
through degraded localization or flight-state estimates.  The plant changes, but the
hidden gap does not.

The key point is that OASG is not a synonym for forecasting error.  A forecast can be
imperfect while the action remains safe, and a forecast can look adequate while the
action becomes unsafe because the observation surface itself is no longer truthful enough
to support legality checks.  ORIUS therefore centers the gap between \emph{what the
controller is allowed to believe} and \emph{what the plant can actually be}.  That gap
is a systems object, not a single-model pathology.

\section{Observed legality versus physical legality}
The intuition can be written compactly even before the full theorem machinery appears.
Let $\hat{x}_t$ denote the observed state used by the nominal controller and let
$x_t^\star$ denote the physical state that is actually realized or could plausibly be
realized given degraded telemetry.  Let $\mathcal{A}(\cdot)$ denote the action set that
is considered admissible.  The OASG event is the regime in which a candidate action
satisfies
\[
u_t \in \mathcal{A}(\hat{x}_t)
\qquad \text{but} \qquad
u_t \notin \mathcal{A}(x_t^\star),
\]
or, equivalently, appears safe on observation while violating the safety predicate on the
true or observation-consistent state.

This notation is intentionally light.  The formal chapters later introduce the precise
uncertainty-set and repair semantics.  At the level of claim boundary, the important
point is simply this: once legality depends on a state that may be stale, delayed,
dropped, or manipulated, the safety argument has to move from nominal action selection to
explicit runtime mediation.

\section{Why ORIUS turns the gap into a runtime object}
That runtime mediation is the central design choice of the book.  ORIUS does not say that
every domain should solve the observability problem with a single estimator, a single
robust controller, or a single proof technique.  It says that each domain must expose the
same sequence of safety-relevant objects at runtime: a reliability judgment on the
observation channel, an uncertainty set for physically plausible state, an admissible
action set under that uncertainty, a repaired or fallback action, and a certificate that
records why the step was allowed.

Once those objects are explicit, the universal and domain-specific parts of the argument
can finally be separated cleanly.  The universal part is the runtime grammar.  The
domain-specific part is how each adapter instantiates state, constraints, repair, and
fallback for its own plant.

\section{Claim classes and promotion discipline}
The claim boundary of the book follows directly from that framing.  ORIUS does not claim
universal optimality, universal hardware closure, or equal empirical maturity across all
domains.  It claims a universal \emph{architecture} and a universal \emph{runtime
contract}.  Evidence is then tiered by what each domain actually clears: witness-depth
reference, defended bounded peer, portability-only row, or experimental boundary row.
This separation is not rhetorical caution; it is the mechanism that
makes a universal monograph defensible under serious review.

The promotion language in the monograph is therefore deliberately typed.  A
\emph{reference witness} carries the deepest theorem-to-code-to-artifact chain and sets
the scientific calibration surface for the rest of the book.  A \emph{defended bounded
peer} clears a domain-specific replay, soundness, fallback, and runtime gate without
inheriting witness-depth proof language automatically.  A \emph{portability-only row}
shows that the architecture can travel structurally but does not yet clear a defended
empirical gate.  An \emph{experimental boundary row} is useful because it marks the edge
of the current program rather than being hidden behind optimistic prose.

\section{Architecture language versus evidence language}
The book therefore uses two kinds of language and keeps them separate.  Architecture
language is universal: true-state violation, constraint margin, repair, certificate,
degraded observation, and intervention.  Evidence language is domain-specific and
promotion-sensitive: some rows close the measured violation gap, some only expose the
structure, and some remain bounded by the current adapter or telemetry surface.  The
monograph is strongest when it refuses to let those two vocabularies collapse into one.

This is also why the universal-first framing is stronger than a witness-first exposition at
the level of exposition.  The reader should learn the common hazard and common runtime
contract first, and only then see how different domains clear different parts of the
evidence ladder.  Battery remains the deepest witness, but it should not monopolize the
conceptual story.

\section{What the book does not claim}
The negative boundary is as important as the positive one.  ORIUS does not claim
universal optimality, universal plant closure, or a flat statement that every domain row
has already reached equal maturity.  It does not claim that uncertainty calibration by
itself proves safety, that runtime certificates replace regulation, or that adapter
portability is equivalent to empirical closure.  The book instead makes a narrower but
more durable claim: one universal safety-layer grammar can govern multiple physical-AI
domains, while a visible parity gate determines how far each domain may currently be
promoted.

The rest of the monograph is written under that contract.  The architecture chapters
explain the reusable runtime object, the theory chapters explain the safety logic under
assumptions, the domain chapters instantiate the contract six times, and the synthesis
chapters make the asymmetry in evidence explicit instead of hiding it.
